% ---- pg-594 （书页 582） ----

(c) 完成最后两个步骤, 并画出显示发生合并时 ESS 值的树图.

\noindent 12.11\quad 假定我们对四个项目 $A,B,C$ 和 $D$ 测量两个变量 $X_1$ 和 $X_2$. 数据如下:

\begin{mvtable}
\setlength{\tabcolsep}{6pt}%
\begin{tabular}{c|cc}
\toprule
& \multicolumn{2}{c}{观测值} \\
\cmidrule(lr){2-3}
\textbf{项 目} & $x_1$ & $x_2$ \\
\midrule
$A$ & 5 & 4 \\
$B$ & 1 & $-2$ \\
$C$ & $-1$ & 1 \\
$D$ & 3 & 1 \\
\bottomrule
\end{tabular}
\end{mvtable}

用 $K$ 均值聚类方法将这四个项目分成 $K=2$ 个聚类. 从初始分组 $(AB)$ 和 $(CD)$ 开始.

\noindent 12.12\quad 重复例 12.11, 从初始分组 $(AC)$ 和 $(BD)$ 开始. 将你得到的结果与例中所给的结果比较. 它们是否相同? 将项目的坐标 $(x_1,x_2)$ 作成点图, 并对解作出评论.

\noindent 12.13\quad 重复例 12.11, 但从项目表的底部开始, 依 $D,C,B,A$ 的顺序进行. 由初始分组 $(AB)$ 和 $(CD)$ 开始. [ 首先可能的重组根据距离 $d^2(D,(AB))$ 和 $d^2(D,(CD))$ 进行. ] 将你的解与例中所给的解进行比较. 它们是否相同? 它们是否应该相同?

\mvsub{下面的练习需要用计算机.}

\noindent 12.14\quad \hyperref[tab:11.9]{表 11.9} 给出 43 种早餐方便粥的 8 个变量的测量值.

\noindent (a) 用表中的数据计算每对方便粥品牌间的欧氏距离.

\noindent (b) 将 (a) 中计算的距离看成是品牌间的相异性, 用单连接和完全连接分层聚类方法对方便粥品牌进行聚类. 构造树图并比较结果.

\noindent 12.15\quad 将\hyperref[tab:11.9]{表 11.9} 中的数据输入一个 $K$ 均值聚类程序. 将方便粥品牌聚类成 $K=2,3,4$ 组. 将结果与练习 12.14 中的结果进行比较.

\noindent 12.16\quad 女子田径运动的国家和地区记录列于\hyperref[tab:1.9]{表 1.9}.

\noindent (a) 用\hyperref[tab:1.9]{表 1.9} 中的数据计算每对国家和地区间的欧氏距离.

\noindent (b) 将距离看成相异性, 用单连接和完全连接分层方法将国家和地区聚类. 构造树图并比较结果.

\noindent (c) 将\hyperref[tab:1.9]{表 1.9} 中的数据输入一个 $K$ 均值聚类程序. 用几个不同的 $K$ 值对国家和地区聚类. 将结果与 (b) 中结果进行比较.

\noindent 12.17\quad 用\hyperref[tab:8.6]{表 8.6} 的男子田径国家和地区记录数据重复练习 12.16. 将结果与练习 12.16 进行比较. 解释任何存在的差别.

\noindent 12.18\quad \hyperref[tab:12.12]{表 12.12} 给出威斯康星及邻近州的 12 个城市间的道路距离. 用多维标度变换法确定这些城市在 $q=1,2,3$ 维空间中的位置. 将最小应力$(q)$对 $q$ 作图并对此图作出解释. 将这些城市的二维多维标度变换结构与一张地图上的城市位置进行比较.

\noindent 12.19\quad \hyperref[tab:12.13]{表 12.13} 给出某些不同年代的考古场所之间的“距离”. 确定这些距离的依据是在这些考古场所发现不同类型陶器的频数. 给定这些距离, 用多维标度变换法确定这些考古场所在 $q=3,4,5$ 维中的坐标. 将最小应力$(q)$对 $q$ 作图并对此图作出解释. 如果可能, 利用 $q=5$ 维解(将场所视作变量)的坐标确定这些场所在二维中的位置(前两个主成分). 注意与场所相联系的年代, 解释二维解的点结构.
